\documentclass[10pt]{article}
\usepackage[margin=0.76in]{geometry}
\usepackage[T1]{fontenc}
\usepackage{lmodern}
\usepackage{amsmath,amssymb,amsthm}
\usepackage{booktabs}
\usepackage{microtype}
\usepackage[hidelinks]{hyperref}
\newtheorem{theorem}{Theorem}
\newtheorem{proposition}{Proposition}
\newcommand{\cO}{\mathcal O}
\newcommand{\cube}{\{0,1\}^{16}}
\title{The Effective-Orbit One-Bit Flip Chain of a Frozen H\'enon Support Core}
\author{Anonymous}
\date{}
\begin{document}
\maketitle

\begin{abstract}
We quotient the sixteen-dimensional support cube by the faithful finite label
action of order $1920$.  The resulting one-bit flip chain has $3024$ states,
constant row sum $16$, and $30240$ nonzero directed entries.  We certify
strong lumpability at every support, orbit-size detailed balance, the exact
flow between four repair-distance levels, and the complete spectrum on
orbit-constant functions.  Its eigenvalues are $16-2k$, with multiplicity
equal to the number of label-group orbits on $k$-subsets.  The repair-zero
diagonal flow is $445696$, independently reproducing the distance-one
correlation from C82.  The result is finite and combinatorial under the
firewall \texttt{NO\_BAD\_EULER\_OR\_ROOT\_NUMBER}.
\end{abstract}

\section{Finite action and quotient}
Let $L$ be the frozen set of sixteen named labels and identify a retained
support $A\subseteq L$ with a vertex of $X=\cube$.  C75 supplies an ambient
lifted action of order $11520$, whose kernel on labels has order $6$.
Accordingly, every support calculation here uses the faithful image
$E\leq\mathfrak S_L$ of order $1920$, never the ambient order as though it
were faithful.  The action has $3024$ orbits on $X$.

Write $\cO_i$ for these support orbits.  For $A\in\cO_i$ define
\[
 q_{ij}=\#\{\ell\in L:A\mathbin{\triangle}\{\ell\}\in\cO_j\}.
\]
Thus $Q=(q_{ij})$ records how the ordinary cube adjacency operator acts on
functions that are constant on $E$-orbits.

\begin{theorem}[Exact strong quotient]
The number $q_{ij}$ is independent of the chosen $A\in\cO_i$.  Every row of
$Q$ sums to $16$, and
\[
 |\cO_i|q_{ij}=|\cO_j|q_{ji}.
\]
Consequently the orbit-size measure is reversible.  The quotient has
$30240$ nonzero directed entries and $15120$ unoriented orbit pairs.
\end{theorem}
\begin{proof}
For $g\in E$, toggling is equivariant:
$g(A\mathbin{\triangle}\{\ell\})=gA\mathbin{\triangle}\{g\ell\}$.
Since $g$ permutes the sixteen labels, it bijects the toggles from any two
representatives of $\cO_i$ and preserves target orbits.  This proves strong
lumpability and the row sum.  Counting actual cube edges from $\cO_i$ to
$\cO_j$ in the two orientations gives the detailed-balance identity.
Support-size parity changes under every toggle, so there are no self-loops;
hence directed entries occur in opposite pairs.
\end{proof}

The number of distinct target orbits in a row has distribution
\[
\begin{array}{c|rrrrrrr}
\toprule
\text{targets}&7&8&9&10&11&12&13\\
\text{rows}&128&384&480&800&864&336&32\\
\bottomrule
\end{array}
\]
and positive entries $1,2,3,4,5$ occur respectively
$19296,7344,1008,1584,1008$ times.

\section{Complete invariant spectrum}
For $B\subseteq L$, let
$\chi_B(A)=(-1)^{|A\cap B|}$ be the Walsh character.  A direct cube
calculation gives adjacency eigenvalue $16-2|B|$.

\begin{theorem}[Invariant Walsh spectrum]
The complete spectrum of $Q$ is $16-2k$, $0\leq k\leq16$.  Its
multiplicity is the number of $E$-orbits on $k$-subsets, namely
\[
\begin{array}{c|rrrrrrrrr}
\toprule
k&0&1&2&3&4&5&6&7&8\\
m_k&1&7&27&73&151&252&352&424&450\\
\midrule
k&9&10&11&12&13&14&15&16\\
m_k&424&352&252&151&73&27&7&1\\
\bottomrule
\end{array}
\]
The multiplicities sum to $3024$; the first and second spectral moments are
$0$ and $77760$.
\end{theorem}
\begin{proof}
The Walsh characters form an orthogonal eigenbasis of functions on $X$.
The action of $E$ permutes them by $g\chi_B=\chi_{gB}$.  Summing characters
over each orbit of subsets produces a nonzero invariant eigenfunction.
Distinct orbit sums have disjoint Walsh support and are independent; every
invariant function has coefficients constant on those orbits.  They
therefore form a basis of the invariant subspace, which is identified with
functions on the $3024$ support orbits.  The stated multiplicities follow.
\end{proof}

\section{Repair flow and predecessor identity}
Let $\rho(A)\in\{0,1,2,3\}$ be the frozen repair distance from C78.  The
numbers of support orbits at these four levels are $1332,1500,180,12$.
Counting actual directed cube edges by the pair of repair levels gives
\[
\begin{array}{c|rrrr}
\toprule
&0&1&2&3\\\midrule
0&445696&40704&0&0\\
1&40704&469376&13184&0\\
2&0&13184&24064&640\\
3&0&0&640&384\\
\bottomrule
\end{array}
\]
whose entries sum to $16\cdot2^{16}$.  In particular, the $(0,0)$ entry is
the number of ordered Hamming-distance-one pairs for which both endpoints
have the full core.  Its value $445696$ exactly equals the C82
distance-one autocorrelation, providing a cross-paper identity obtained by
a different aggregation route.

\section{Certificate and scope}
The producer explicitly generates all $1920$ label permutations, whereas
the independent checker constructs support-orbit components using only five
named generators.  Both verify every one of the $65536$ supports and every
quotient row.  A symbolic radial-cube check, clean replay, and twenty hostile
mutations complete the audit.  The canonical evidence SHA-256 is
\begin{center}
\texttt{7b3e2179590c3dc8662a59f1d79ffbb1}\\[-2pt]
\texttt{2f2a4a787438a6902d6c28b2842e70b8}.
\end{center}
No arithmetic/local data, Euler factors, root numbers, automorphy, full
Burnside ring or table of marks, or Hilbert--Polya operator is claimed.
\end{document}
